\section{Local Smith factors of the full exterior tensor}
\label{sec:fock-smith}

The coefficient-rank question can also be examined before imposing
particle and energy restrictions. In that larger space the local
valuation spectrum is explicit. The restrictions remain essential:
their Smith factors need not agree with those of the larger tensor.

Let $\kk$ have characteristic two and put $R=\kk[[X]]$.
For an $M$-mode space, its full exterior space is
$\bigoplus_{r=0}^M\bigwedge^r\kk^M$, with basis indexed by all
subsets of the modes. Let $p$ be an involution of the mode indices,
$P$ its permutation matrix, and $w_i\geq0$ integers. Suppose
$C\in\operatorname{GL}_M(R)$, and set
\[
 J=P\operatorname{diag}(X^{w_i}),\qquad
 K=\begin{pmatrix}C&XCJ\\XCJ&C\end{pmatrix}.
\]
Define a four-party tensor $\mathcal F_K(A,B,C_0,D)$ by the minor
with row sets $(A,C_0)$ and column sets $(D,B)$ in these two blocks.
Its value is zero when the total row and column cardinalities differ.
All four terminals use the full exterior space.

\begin{proposition}[The unrestricted exterior valuation spectrum]
\label{prop:fock-smith-spectrum}
The Smith valuations of the $ad\mid bc$ flattening of $\mathcal F_K$
over $R$ are the multiset
\[
 \left\{\sum_{i\in U}(1+w_i)+
              \sum_{i\in V}(1+w_{p(i)}):
              U,V\subseteq[M]\right\}.
\]
Its multiplicity generating polynomial, with ordinary integer
coefficients, is
\[
 \prod_{i=1}^M(1+t^{1+w_i})(1+t^{1+w_{p(i)}}).
\]
Tensoring two such tensors and grouping their four parties adds
the valuation multisets. For identical mode-weight multisets this
gives $\prod_i(1+t^{1+w_i})^4$.

Multiplication or division by a scalar formal unit does not change
these valuations. No equality of spectra is asserted after
restricting the terminal particle or energy classes.
\end{proposition}
\begin{proof}
Multiply $K$ on the left by
$\operatorname{diag}(C^{-1},C^{-1})$. On the full exterior tensor
this applies separate invertible exterior maps at terminals $a,c$.
Their matrices and inverses have entries in $R$, so they preserve
the Smith valuations of the adjacent flattening. Next permute the
second row block and the second column block by $P$. Since
$P^2=I$, the resulting one-particle matrix is a direct sum over
modes of
\[
 \begin{pmatrix}1&t_i\\s_i&1\end{pmatrix},\qquad
 s_i=X^{1+w_i},\quad t_i=X^{1+w_{p(i)}}.
\]
All minors consequently factor over the modes. Characteristic two
removes the signs from reordering exterior factors.

For one mode, list row indices $(a,d)$ and column indices $(b,c)$
as $00,01,10,11$. The full-minor flattening is
\[
 \begin{pmatrix}
 1&0&0&1\\
 0&s_i&0&0\\
 0&0&t_i&0\\
 1&0&0&1+s_it_i
 \end{pmatrix}.
\]
Adding its first row to its last row and then its first column to
its last column makes it diagonal, with entries
$1,s_i,t_i,s_it_i$. These are invertible row and column operations
over $R$. Tensoring their diagonal forms proves the stated
multiset; it also proves the tensor-product assertion. The
multiplicities count distinct diagonal positions as integers and
must not be reduced modulo two. Finally a scalar unit is itself
an invertible row operation. A terminal restriction instead takes
a submatrix, for which this equality of Smith spectra need not
persist.
\end{proof}

For the finite digit kernel,
$w_i=2(q-1)E_q(i)$ and $p(i)=M-1-i$. Its matrix $C$ is a tensor
product of $\left(\begin{smallmatrix}1&1\\1&X^{2(q-1)q^j}
\end{smallmatrix}\right)$, each invertible at zero. Thus the
proposition gives two unrestricted exterior species with mode
energies $\ell_q(i)=1+w_i$ for one minor, and four for the product
of the two minors. The common divisor removed at $X=1$ has value
one at zero and hence does not alter this local unrestricted
spectrum. A better estimate must use the selected classes or
additional cancellation between coefficients.

\subsection{Exact local spectra in two finite restricted classes}

The restricted matrices can be substantially smaller than the
full exterior spaces. For radix four, size-two symbols of combined
energy $64$, the dimension-eight highest-normalizer coefficient
has the following $64$ Smith valuations across $ad\mid bc$:
\[
\begin{aligned}
 &0^1,\ 38^6,\ 56^4,\ 64^{10},\ 70^2,\ 72^2,\ 76^2,\ 128^{10},\\
 &180^2,\ 184^2,\ 186^2,\ 192^{10},\ 200^4,\ 218^6,\ 256^1.
\end{aligned}
\]
Here a superscript is a multiplicity. All factors are resolved by
exact elimination modulo $X^{257}$. As a separate finite check,
the matrix of multiplication modulo $X^{97}$ has binary rank
$1091$, equal to $\sum_v\max(97-v,0)$ over this multiset. The
linear dimension-two control has valuations $0^1,8^2,16^1$.

For a single normalized middle minor in eight modes, restrict each
terminal to the four-element subsets of digit energy $42$.
There are eight such subsets, and the normalized output is a
dimension-eight crossing. The common divisor has degree $248$;
the divided polynomial has degree $768$ and scalar coefficient
span dimension $51$. Its coordinates have no additional common
polynomial factor. The best contiguous low/high coefficient
grouping has ranks $1$ and $28$. Its Smith spectrum has $27$
factors of valuation at most the midpoint $384$; the preceding
size-two-symbol example also has $27$ at its midpoint $96$.

These exact local calculations do not optimize arbitrary
decompositions. In particular a Smith-state count controls a
valuation-based factorization, while sums of coefficient matrices
can still cancel. The finite results leave that cancellation
question open.
